\documentclass[11pt]{article}

\usepackage[margin=1in]{geometry}
\usepackage[T1]{fontenc}
\usepackage{lmodern}
\usepackage{amsmath,amssymb,amsthm}
\usepackage{mathtools}
\numberwithin{equation}{section}
\usepackage{bm}
\usepackage{booktabs}
\usepackage{tabularx}
\usepackage{array}
\usepackage{graphicx}
\usepackage{siunitx}
\usepackage{xcolor}
\usepackage{enumitem}
\usepackage{microtype}
\usepackage{algorithm}
\usepackage{algpseudocode}
\usepackage{tikz}
\usetikzlibrary{arrows.meta,positioning,calc}
\usepackage[round,numbers,sort&compress]{natbib}
\usepackage[colorlinks=true,linkcolor=blue,citecolor=blue,urlcolor=blue]{hyperref}

% ---------------------------------------------------------------- shortcuts
\newcommand{\R}{\mathbb{R}}
\newcommand{\E}{\mathbb{E}}
\newcommand{\Prob}{\mathbb{P}}
\newcommand{\dd}{\,\mathrm{d}}
\newcommand{\pdd}[2]{\frac{\partial #1}{\partial #2}}
\newcommand{\pddt}[2]{\frac{\partial^{2} #1}{\partial #2^{2}}}
\newcommand{\norm}[1]{\left\lVert #1 \right\rVert}
\newcommand{\abs}[1]{\left\lvert #1 \right\rvert}
\newcommand{\xx}{\mathbf{x}}
\newcommand{\qq}{\mathbf{q}}
\newcommand{\ww}{\mathbf{w}}
\newcommand{\vv}{\mathbf{v}}
\newcommand{\DD}{\mathbf{D}}
\newcommand{\T}{\mathsf{T}}
\newcommand{\CVaR}{\operatorname{CVaR}}
\newcommand{\VaR}{\operatorname{VaR}}
\newcommand{\argmin}{\operatorname*{arg\,min}}
\newcommand{\argmax}{\operatorname*{arg\,max}}
\newcommand{\ind}{\mathbf{1}}
\newcommand{\me}{\mathrm{e}}
\newcommand{\Proj}{\operatorname{Proj}}
\newcommand{\tabhead}[1]{\textbf{#1}}
\newcolumntype{Y}{>{\raggedright\arraybackslash}X}

\theoremstyle{plain}
\newtheorem{proposition}{Proposition}
\newtheorem{remark}{Remark}
\newtheorem{assumption}{Assumption}

\title{Algorithms for the Stochastic Pump-and-Treat Remediation Problem\\[2pt]
\large Sample-average approximation, Rockafellar--Wets progressive hedging,
stochastic gradients, decomposition, and surrogate methods}
\author{Groundwater Remediation Models}
\date{\today}

\begin{document}
\maketitle

\begin{abstract}
This document surveys and compares algorithms for the PDE-constrained,
chance-constrained stochastic optimization problem formulated in
\texttt{optimization.tex}.  The decision is a here-and-now vector of extraction
rates; the uncertainty is a random conductivity field and, optionally, random
source, boundary, and transport parameters; the objective is an expected
life-cycle cost; and the constraints include probabilistic compliance and
drawdown requirements.  We describe the sample-average approximation (SAA)
baseline, the Rockafellar--Wets progressive hedging algorithm (PHA) and its
use with chance constraints, the L-shaped/Benders family, stochastic-gradient
and stochastic-approximation methods, derivative-free and global search,
surrogate and reduced-order optimization, and chance-constraint-specific
reformulations.  The emphasis is on \emph{which algorithm to use when}, on the
information each method requires (gradients, scenario decomposability,
parallelism), and on the verification checks that keep a stochastic
optimization run honest.  A hybrid roadmap is given at the end: coarse
derivative-free exploration, adjoint-based SAA refinement, PHA or a surrogate
for large scenario sets, and an independent out-of-sample certificate.
\end{abstract}

\tableofcontents

% ======================================================================
\section{Purpose and scope}
\label{sec:purpose}

The forward-model chain (one-, two-, and three-dimensional flow and transport,
the Ogata--Banks reference, and the Fourier neural operator) is documented in
\texttt{formulation.tex}, \texttt{2D\_plan.md}, \texttt{3D\_plan.md}, and
\texttt{optimization.tex}.  This document is the algorithm layer: it assumes
that a trustworthy scenario generator and a forward map
$(\mathbf{h},\mathbf{C})=\mathcal{S}(\qq,\omega)$ exist, and it asks how to
solve the resulting stochastic program.

The problem has four features that determine the algorithm choice:
\begin{enumerate}[leftmargin=1.6em,itemsep=2pt]
  \item \textbf{PDE constraints.}  Every objective and constraint evaluation
        requires a flow and transport solve.  The state is high-dimensional,
        but the decision $\qq$ is low- or moderate-dimensional
        ($N_q=N_wN_p$, typically $10$--$100$).
  \item \textbf{Random fields.}  The conductivity is a log-normal random field
        with a Karhunen--Lo\`eve expansion.  The scenario space is high-
        dimensional, but the cost and constraints depend on it through a
        relatively small number of decision-relevant functionals.
  \item \textbf{Chance constraints.}  The compliance and drawdown
        requirements are probabilistic.  Their empirical counterparts are
        step functions of $\qq$, hence nonsmooth and nonconvex.
  \item \textbf{Expensive scenarios.}  A single high-fidelity scenario solve
        can take seconds to minutes; the final reliability certificate may
        need thousands of scenarios.  Parallelism over scenarios is available
        and should be exploited.
\end{enumerate}

The algorithms below differ mainly in how they handle these four features.
No single method is best for all regimes; the practical answer is a hybrid,
with each stage chosen to match the information and budget available.

% ======================================================================
\section{The stochastic program in compact form}
\label{sec:problem}

\subsection{Random inputs and scenarios}

Let $(\Omega,\mathcal{F},\Prob)$ be a probability space and let
$\omega\in\Omega$ index a scenario.  The baseline random field is
\begin{equation}
\label{eq:lognormal}
K(\xx,\omega)=K_0\exp\bigl(G(\xx,\omega)\bigr),
\qquad
G(\xx,\omega)=\sum_{j=1}^{N_{KL}}\sqrt{\lambda_j}\,\phi_j(\xx)\,\xi_j(\omega),
\end{equation}
with $\xi_j\sim\mathcal{N}(0,1)$ i.i.d., and $(\lambda_j,\phi_j)$ the
eigenpairs of the covariance operator.  The truncation $N_{KL}$ is chosen to
retain $95$--$99\%$ of the variance.  Optional random inputs (porosity, source
strength and location, boundary head, inlet concentration, dispersivities) are
appended to $\omega$ in the same way.  A scenario is the vector of all random
variables needed to run one forward solve; it is stored as coefficients, not
only as a realized field.

\subsection{Decision variables and feasible set}

For a steady policy the decision is
\begin{equation}
\label{eq:qsteady}
\qq=(q_1,\dots,q_{N_w})\in\R^{N_w},
\qquad q_i\ge 0,
\end{equation}
and for a pumping schedule it is
\begin{equation}
\label{eq:qpiecewise}
\qq=(q_{ik})_{i=1,\dots,N_w}^{k=1,\dots,N_p}\in\R^{N_wN_p},
\end{equation}
with $q_i(t)=\sum_k q_{ik}\ind_{[t_{k-1},t_k)}(t)$.  The admissible set is the
convex polytope
\begin{equation}
\label{eq:Qad}
\mathcal{Q}_{\mathrm{ad}}
=\left\{\qq:\
\begin{aligned}
&0\le q_{ik}\le q_i^{\max},\\
&\sum_{i=1}^{N_w}q_{ik}\le Q_{\max}, && k=1,\dots,N_p,\\
&\sum_{i=1}^{N_w}\sum_{k=1}^{N_p}q_{ik}\Delta t_k\le V_{\max},
\end{aligned}
\right\}.
\end{equation}
The difficulty is not the rate bounds; it is the PDE state map and the
probability constraints.  We write $N_q=N_wN_p$ for the number of decision
variables.

\subsection{Cost, violations, and the full program}

For a scenario $\omega$ and a policy $\qq$, over a horizon of $T$ years the
objective is the sum of pumping costs and terminal remediation penalties:
\begin{equation}
\label{eq:cost}
J(\qq,\omega)=
\underbrace{\sum_{i=1}^{N_w}\sum_{k=1}^{N_p}
\bigl(c_e+c_w\,C(\xx_i,t_k;\qq,\omega)\bigr)q_{ik}\Delta t_k}
_{\text{pumping and treatment cost}}
+
\underbrace{c_m\int_\Omega\phi\,C(\xx,T;\qq,\omega)\dd\xx
+c_\tau\min(\tau_c,T)
+c_{\mathrm{fail}}\ind\{\tau_c>T\}}
_{\text{remediation penalty after }T}.
\end{equation}
Here $c_e$ is the energy cost per unit pumped volume, $c_w$ the treatment cost
per unit extracted concentration, $c_m$ the terminal residual-mass penalty,
$c_\tau$ the cleanup-time penalty, and $c_{\mathrm{fail}}$ the penalty for
missing the cleanup target.  The compliance and drawdown violations are
\begin{align}
G(\qq,\omega)
&=\max_{\xx_c\in\mathcal{M}_c}\ \max_{t\in[T_c,T]}
\bigl(C(\xx_c,t;\qq,\omega)-C_{\mathrm{reg}}\bigr),
\label{eq:G}\\
H(\qq,\omega)
&=\min_{\xx_o\in\mathcal{O}_h}\ \min_{t\in(0,T]}
\bigl(h(\xx_o,t;\qq,\omega)-h_{\min}(\xx_o)\bigr),
\label{eq:H}
\end{align}
so $G\le0$ means compliance and $H\ge0$ means acceptable drawdown.  The full
stochastic program is
\begin{equation}
\label{eq:stochprogram}
\boxed{
\begin{aligned}
\min_{\qq\in\mathcal{Q}_{\mathrm{ad}}}\quad
& \E\bigl[J(\qq,\omega)\bigr]\\
\text{s.t.}\quad
& \Prob\bigl(G(\qq,\omega)\le0\bigr)\ge1-\alpha,\\
& \Prob\bigl(H(\qq,\omega)\ge0\bigr)\ge1-\alpha_h,\\
& (\mathbf{h},\mathbf{C})(\cdot,\cdot;\qq,\omega)
  =\mathcal{S}(\qq,\omega)
  \quad\text{a.s. }\omega\in\Omega.
\end{aligned}}
\end{equation}
The decision is here-and-now: it is chosen before $\omega$ is observed and is
not adapted to the realized conductivity field.

\subsection{Sample-average approximation}

Draw $M$ independent scenarios $\omega^1,\dots,\omega^M$.  The SAA program is
\begin{equation}
\label{eq:saa}
\begin{aligned}
\min_{\qq\in\mathcal{Q}_{\mathrm{ad}}}\quad
& J_M(\qq)=\frac{1}{M}\sum_{m=1}^{M}J(\qq,\omega^m)\\
\text{s.t.}\quad
& \frac{1}{M}\sum_{m=1}^{M}\ind\{G(\qq,\omega^m)\le0\}\ge1-\alpha,\\
& \frac{1}{M}\sum_{m=1}^{M}\ind\{H(\qq,\omega^m)\ge0\}\ge1-\alpha_h.
\end{aligned}
\end{equation}
The empirical probabilities are step functions, so \eqref{eq:saa} is a
nonsmooth, nonconvex program.  Three tractable reformulations are used below:
the exact mixed-integer form, the CVaR surrogate, and the sigmoid smoothing.
The SAA is not the only possible approximation (scenario reduction,
importance sampling, and multilevel Monte Carlo also apply), but it is the
reference because it separates the statistical error from the optimization
error and supports out-of-sample validation.

\subsection{Structural properties that drive the algorithm choice}

\begin{table}[htbp]
\centering
\small
\begin{tabularx}{\textwidth}{@{}lY@{}}
\toprule
\tabhead{Property} & \tabhead{Consequence for algorithms}\\
\midrule
$N_q=N_wN_p\approx10$--$100$ & Reduced-space methods (adjoint, derivative-free,
surrogate) are natural; full-space Newton methods are possible but larger.\\
$M\approx10^2$--$10^4$ scenarios & Scenario decomposition and parallel evaluation
matter; deterministic-equivalent NLP is practical only for moderate $M$.\\
PDE state is high-dimensional & Adjoint gradients cost about one forward solve;
automatic differentiation through the solver is an alternative.\\
Cost is smooth in $\qq$ (if the indicator is removed) & L-BFGS-B, SQP, trust
region, and stochastic gradient apply.\\
Chance constraint is nonsmooth & Use MINLP, CVaR, sigmoid smoothing, or
derivative-free search; do not feed a raw indicator to a gradient method.\\
Objective and PDE are scenario-separable & PHA, dual decomposition, and
scenario-parallel SAA apply.\\
Chance constraint couples scenarios & Needs a dual/penalty layer, CVaR with an
auxiliary variable, or a global (non-decomposed) solve.\\
Forward solve is expensive & Surrogate, multi-fidelity, Bayesian optimization,
and active learning become attractive.\\
\bottomrule
\end{tabularx}
\caption{Structural properties of the remediation problem and their algorithmic
consequences.}
\label{tab:structure}
\end{table}

% ======================================================================
\section{Algorithm selection criteria}
\label{sec:selection}

Table~\ref{tab:methods} is the one-page summary.  The columns are the
information the method needs, the structure it exploits, and the regime in
which it is the recommended choice.

\begin{table}[htbp]
\centering
\scriptsize
\setlength{\tabcolsep}{3pt}
\begin{tabularx}{\textwidth}{@{}>{\raggedright\arraybackslash}p{2.9cm}cccc>{\raggedright\arraybackslash}X@{}}
\toprule
\tabhead{Method} & \tabhead{Grad.} & \tabhead{Decomp.} &
\tabhead{Chance} & \tabhead{Par.} & \tabhead{Use / caveat}\\
\midrule
SAA + adjoint + L-BFGS-B/SQP & exact & no & CVaR/sigmoid & evals & Baseline for
smooth $N_q\lesssim100$, moderate $M$; local.\\
Rockafellar--Wets PHA & subproblem & yes & CVaR/Lagrangian & strong & Large
$M$, separable subproblems; convex guarantee, heuristic otherwise.\\
L-shaped / Benders & cuts & master--sub & limited & yes & Linear/quadratic
recourse; generalized Benders for smooth nonlinear.\\
Stochastic gradient / Adam & stochastic & mini-batch & sigmoid/CVaR & strong &
Many scenarios, online; noisy, step-size tuning.\\
Pattern search / MADS & none & no & exact empirical & yes & Small $N_q$,
nonsmooth; local.\\
CMA-ES / differential evolution & none & no & exact empirical & yes & Small
$N_q$, rugged; expensive, no guarantee.\\
Bayesian optimization & none & no & constrained EI & no & Very expensive FV,
$N_q\lesssim10$; hard for many constraints.\\
FNO/POD surrogate + trust region & surrogate & yes & smooth & training &
Many scenarios; surrogate bias must be validated.\\
Robust / DRO & subgradient & no & robust & limited & Distribution shift;
conservative; needs an ambiguity set.\\
\bottomrule
\end{tabularx}
\caption{Outer optimization methods for the stochastic remediation problem.}
\label{tab:methods}
\end{table}

The rest of the document discusses each family in more detail.  Section~\ref{sec:saa-nlp}
covers the SAA/NLP baseline, Section~\ref{sec:pha} the Rockafellar--Wets PHA,
Section~\ref{sec:decomp} other decompositions, Section~\ref{sec:sg} stochastic
gradients, Section~\ref{sec:df} derivative-free methods, Section~\ref{sec:surrogate}
surrogates, and Section~\ref{sec:chance} chance-constraint-specific algorithms.
Section~\ref{sec:hybrid} gives the recommended hybrid.

% ======================================================================
\section{Sample-average approximation with a deterministic NLP solver}
\label{sec:saa-nlp}

\subsection{Reduced-space formulation}

For a fixed scenario $\omega^m$, let $J_m(\qq)=J(\qq,\omega^m)$ and
$G_m(\qq)=G(\qq,\omega^m)$.  After choosing a smooth reformulation of the
chance constraint, the SAA problem has the form
\begin{equation}
\label{eq:saa-nlp}
\min_{\qq\in\mathcal{Q}_{\mathrm{ad}}}\quad
J_M(\qq)=\frac{1}{M}\sum_{m=1}^{M}J_m(\qq)
\quad\text{s.t.}\quad
c_M(\qq)\le0,
\end{equation}
where $c_M$ is, for example, the CVaR constraint
\begin{equation}
\label{eq:cvar-constraint}
c_M(\qq)=\min_{\eta\in\R}\left\{
\eta+\frac{1}{\alpha M}\sum_{m=1}^{M}
\bigl(G_m(\qq)-\eta\bigr)_+
\right\}\le0,
\end{equation}
or the sigmoid constraint
\begin{equation}
\label{eq:sigmoid-constraint}
c_M(\qq)=1-\alpha-\frac{1}{M}\sum_{m=1}^{M}
\sigma\!\left(-\frac{G_m(\qq)}{\varepsilon_s}\right)\le0,
\qquad
\sigma(z)=\frac{1}{1+\me^{-z}}.
\end{equation}
The decision dimension is only $N_q$, so a reduced-space NLP solver is the
natural baseline.  The state variables are eliminated by the forward solve;
the solver sees only $\qq$ and the aggregated cost and constraint.

\subsection{Adjoint gradients}

For a smooth cost $J_m$ and no chance constraint, the reduced gradient is
computed by the discrete adjoint.  Write the discretized transport residual as
$\mathbf{F}_h(\mathbf{C}_h,\qq)=0$, form the Lagrangian
\begin{equation}
\label{eq:discreteL}
\mathcal{L}_h(\mathbf{C}_h,\qq,\bm{\psi}_h)
=J_h(\mathbf{C}_h,\qq)
-\bm{\psi}_h^{\T}\mathbf{F}_h(\mathbf{C}_h,\qq),
\end{equation}
and solve
\begin{equation}
\label{eq:discreteadj}
\left(\frac{\partial\mathbf{F}_h}{\partial\mathbf{C}_h}\right)^{\T}
\bm{\psi}_h
=\left(\frac{\partial J_h}{\partial\mathbf{C}_h}\right)^{\T},
\qquad
\frac{\dd J_h}{\dd\qq}
=\frac{\partial J_h}{\partial\qq}
-\left(\frac{\partial\mathbf{F}_h}{\partial\qq}\right)^{\T}
\bm{\psi}_h.
\end{equation}
The adjoint system is linear, has the sparsity of the linearized transport
operator, and is solved backward in time.  Its cost is comparable to one
forward solve and is independent of $N_q$.  This is why the adjoint is the
recommended gradient for $N_w\gg1$.

For the CVaR constraint \eqref{eq:cvar-constraint}, the gradient with respect
to $\qq$ is
\begin{equation}
\label{eq:cvar-grad}
\nabla_{\qq} c_M(\qq)
=\frac{1}{\alpha M}\sum_{m\in\mathcal{A}(\qq)}
\nabla_{\qq}G_m(\qq),
\end{equation}
where $\mathcal{A}(\qq)=\{m:G_m(\qq)>\eta^\star(\qq)\}$ is the active tail
set and $\eta^\star$ is a minimizer in \eqref{eq:cvar-constraint}.  The
gradient of $G_m$ is again obtained by an adjoint of the transport equation
with the max-over-monitoring-points and max-over-time resolved (or smoothed).
For the sigmoid constraint,
\begin{equation}
\label{eq:sigmoid-grad}
\nabla_{\qq} c_M(\qq)
=\frac{1}{M\varepsilon_s}\sum_{m=1}^{M}
\sigma\!\left(-\frac{G_m}{\varepsilon_s}\right)
\left[1-\sigma\!\left(-\frac{G_m}{\varepsilon_s}\right)\right]
\nabla_{\qq}G_m(\qq).
\end{equation}
Small $\varepsilon_s$ makes the sigmoid close to the indicator but also makes
the gradient small when all scenarios are far from the boundary.  Use a
continuation: start with $\varepsilon_s=10^{-1}$, then $10^{-2}$, then
$10^{-3}$, warm-starting each stage.

\subsection{Solvers and constraint handling}

\begin{itemize}[leftmargin=1.6em,itemsep=2pt]
  \item \textbf{L-BFGS-B} for box constraints and a smooth penalized objective.
        It is the simplest baseline and scales well to $N_q\approx100$.
        Penalize the chance constraint with an augmented Lagrangian or a
        quadratic penalty and increase the penalty if the constraint is
        violated at the candidate solution.
  \item \textbf{SQP or interior-point} for the NLP \eqref{eq:saa-nlp} when
        accurate constraint gradients are available and $N_q$ is small
        (tens).  These methods handle the CVaR epigraph variable $\eta$ and
        the smooth sigmoid constraint directly.
  \item \textbf{Trust-region} methods when the objective is noisy (finite-
        difference gradients, surrogate, or Monte Carlo).  The trust region
        limits the step to the region where the model is reliable.
  \item \textbf{Multi-start} from a space-filling design (Latin hypercube or
        Sobol) because the SAA problem is nonconvex.  Keep the best few
        candidates and refine them.
\end{itemize}

\subsection{Algorithm}

\begin{algorithm}[htbp]
\caption{SAA with adjoint gradients and a smooth chance constraint}
\label{alg:saa}
\begin{algorithmic}[1]
\Require scenario archive $\{\omega^m\}_{m=1}^{M}$, test archive, initial
         $\qq^0$, smoothing schedule $\varepsilon_s$, tolerances
\Ensure candidate policy $\qq^\star$, SAA objective, constraint value
\State Generate and store the training and test archives with independent
       seeds; record $N_{KL}$ and all random coefficients
\State Precompute unit-rate flow responses and store the well operators
\For{each $\varepsilon_s$ in the continuation schedule}
  \State Warm-start from the previous stage (or $\qq^0$ for the first stage)
  \Repeat
    \State For each scenario $m$, solve the forward PDE at $\qq$ and compute
           $J_m$, $G_m$, $H_m$
    \State Compute $\nabla J_M$ and $\nabla c_M$ by the discrete adjoint
           (or by automatic differentiation through the solver)
    \State Take an L-BFGS-B/SQP step on the augmented Lagrangian
           $J_M+\rho\,[c_M]_+^2$ (or on the CVaR epigraph)
    \State Update the penalty $\rho$ if the constraint violation did not
           decrease
  \Until{stationarity and constraint tolerance are met}
\EndFor
\State Evaluate $\qq^\star$ on the test archive; report
       $\hat J_{\mathrm{test}}$, $\hat p_{\mathrm{test}}$,
       $\widehat{\CVaR}_{\mathrm{test}}$
\State Repeat with $K$ independent SAA replications; report the best
       candidate and the spread
\State \Return $\qq^\star$ and the out-of-sample certificate
\end{algorithmic}
\end{algorithm}

\paragraph{Strengths.}  Exact gradients, mature solvers, good local convergence,
and a clear path to a reliability certificate.  This is the baseline against
which every other method should be compared.

\paragraph{Weaknesses.}  The SAA NLP is nonconvex; the raw chance constraint is
nonsmooth; and the method evaluates all $M$ scenarios at every iteration.  For
$M$ in the thousands and an expensive forward solve, the cost is dominated by
the scenario sweeps.

% ======================================================================
\section{Rockafellar--Wets progressive hedging algorithm}
\label{sec:pha}

\subsection{Scenario decomposition and nonanticipativity}

The SAA objective and the PDE state equations are separable across scenarios;
the only coupling is that the same here-and-now decision $\qq$ must be used in
every scenario.  Introduce one copy $\qq^m$ per scenario and the
nonanticipativity (consensus) constraints
\begin{equation}
\label{eq:consensus}
\qq^m=\qq,\qquad m=1,\dots,M.
\end{equation}
Then \eqref{eq:saa} is equivalent to
\begin{equation}
\label{eq:consensus-program}
\min_{\qq,\{\qq^m\}}\ \sum_{m=1}^{M}\pi_m f_m(\qq^m)
\quad\text{s.t.}\quad
\qq^m=\qq,\ \ \qq^m\in\mathcal{Q}_{\mathrm{ad}},
\end{equation}
where $\pi_m=1/M$ (or the scenario probabilities) and $f_m$ is the
scenario-$m$ objective including its local constraints.  Without a chance
constraint, $f_m(\qq^m)=J_m(\qq^m)$.  The chance constraint is handled in
Section~\ref{sec:pha-chance}.

The progressive hedging algorithm of Rockafellar and Wets
\citep{rockafellar1991} is an augmented Lagrangian method for
\eqref{eq:consensus-program}.  It does not solve the deterministic equivalent
as one large NLP.  Instead, it solves $M$ independent scenario subproblems per
iteration, aggregates their solutions, and updates multipliers that penalize
disagreement.  The name ``progressive hedging'' reflects the fact that each
scenario is progressively hedged toward the consensus policy.

\subsection{Augmented Lagrangian derivation}

For a penalty parameter $r>0$ and multiplier vectors $\ww^m\in\R^{N_q}$,
define the augmented Lagrangian
\begin{equation}
\label{eq:pha-L}
L_r(\qq,\{\qq^m\},\{\ww^m\})
=\sum_{m=1}^{M}\pi_m\left[
f_m(\qq^m)
+(\ww^m)^{\T}(\qq^m-\qq)
+\frac{r}{2}\norm{\qq^m-\qq}^2
\right],
\end{equation}
with $\qq\in\mathcal{Q}_{\mathrm{ad}}$ and $\qq^m\in\mathcal{Q}_{\mathrm{ad}}$.
The PHA iteration is
\begin{align}
\qq^{k+1}
&=\argmin_{\qq\in\mathcal{Q}_{\mathrm{ad}}}\ L_r(\qq,\{\qq^{k,m}\},\{\ww^{k,m}\}),
\label{eq:pha-q}\\
\qq^{k+1,m}
&=\argmin_{\qq^m\in\mathcal{Q}_{\mathrm{ad}}}\
\pi_m f_m(\qq^m)+(\ww^{k,m})^{\T}\qq^m
+\frac{r}{2}\norm{\qq^m-\qq^{k+1}}^2,
\label{eq:pha-sub}\\
\ww^{k+1,m}
&=\ww^{k,m}+r\bigl(\qq^{k+1,m}-\qq^{k+1}\bigr).
\label{eq:pha-w}
\end{align}
The aggregate step \eqref{eq:pha-q} has the closed form
\begin{equation}
\label{eq:pha-aggregate}
\qq^{k+1}
=\Proj_{\mathcal{Q}_{\mathrm{ad}}}
\left(
\sum_{m=1}^{M}\pi_m\qq^{k,m}
+\frac{1}{r}\sum_{m=1}^{M}\pi_m\ww^{k,m}
\right).
\end{equation}
If $\ww^{0,m}=0$ for all $m$, the weighted multiplier sum remains zero, and
when $\mathcal{Q}_{\mathrm{ad}}$ is convex and the scenario solutions are
feasible, \eqref{eq:pha-aggregate} reduces to the weighted average
\begin{equation}
\label{eq:pha-average}
\qq^{k+1}=\sum_{m=1}^{M}\pi_m\qq^{k+1,m}.
\end{equation}
This average is the policy actually reported; the individual $\qq^{k,m}$ are
auxiliary scenario policies.

\begin{remark}[Why the proximal term matters]
The quadratic term $(r/2)\norm{\qq^m-\qq}^2$ makes the scenario subproblem
strongly convex near the aggregate and prevents a scenario from moving too far
from the consensus in one iteration.  It also makes the subproblem
differentiable even when $f_m$ is not.  A pure linear penalty (without the
quadratic term) is not a valid augmented Lagrangian for this consensus
problem and can fail to converge.
\end{remark}

\subsection{The PHA algorithm}

\begin{algorithm}[htbp]
\caption{Rockafellar--Wets progressive hedging for the SAA problem}
\label{alg:pha}
\begin{algorithmic}[1]
\Require scenarios $\{\omega^m\}_{m=1}^{M}$, weights $\pi_m$, penalty $r>0$,
         initial policy $\qq^0$, tolerances $\varepsilon_{\mathrm{prim}},
         \varepsilon_{\mathrm{step}}$, max iterations $K_{\max}$
\Ensure consensus policy $\qq^\star$ and scenario multipliers
\State Initialize $\qq^{0,m}=\qq^0$, $\ww^{0,m}=0$ for all $m$
\For{$k=0,1,\dots,K_{\max}-1$}
  \State \textbf{(scenario subproblems, in parallel)} for each $m$:
         \Statex \hspace{1.2em} solve
         $\displaystyle \qq^{k+1,m}=\argmin_{\qq^m\in\mathcal{Q}_{\mathrm{ad}}}
         \ \pi_m f_m(\qq^m)+(\ww^{k,m})^{\T}\qq^m
         +\frac{r}{2}\norm{\qq^m-\qq^{k,m}}^2$
         \Statex \hspace{1.2em} (warm-start from $\qq^{k,m}$; use an adjoint
         gradient and a few L-BFGS-B/trust-region iterations)
  \State \textbf{(aggregation)} compute
         $\displaystyle \qq^{k+1}
         =\Proj_{\mathcal{Q}_{\mathrm{ad}}}
         \Bigl(\sum_m\pi_m\qq^{k+1,m}
         +\frac{1}{r}\sum_m\pi_m\ww^{k,m}\Bigr)$
  \State \textbf{(multiplier update)} for each $m$:
         $\ww^{k+1,m}=\ww^{k,m}+r(\qq^{k+1,m}-\qq^{k+1})$
  \State \textbf{(stopping)} stop if
         $\max_m\norm{\qq^{k+1,m}-\qq^{k+1}}\le\varepsilon_{\mathrm{prim}}$
         and $\norm{\qq^{k+1}-\qq^{k}}\le\varepsilon_{\mathrm{step}}$
         \State Optionally adapt $r$: increase it if the consensus residual
         stalls; decrease it if the multipliers oscillate
\EndFor
\State \Return $\qq^\star=\qq^{K}$, $\{\ww^{K,m}\}$, and the scenario
       objectives $\{f_m(\qq^{K,m})\}$
\end{algorithmic}
\end{algorithm}

\subsection{Convergence and practical behavior}

\begin{proposition}[Convex case]
\label{prop:pha-convex}
If the scenario objectives $f_m$ are convex, $\mathcal{Q}_{\mathrm{ad}}$ is
convex and closed, the Slater condition holds, and the augmented Lagrangian
subproblems are solved exactly, then PHA converges to an optimal solution of
\eqref{eq:consensus-program}; the multipliers converge to optimal
nonanticipativity prices \citep{rockafellar1991}.
\end{proposition}

The nonconvex case is different.  The PDE-constrained cost is not convex in
$\qq$ in general, and the chance constraint is not convex even when the
dynamics are linear.  There is no general global-convergence guarantee.  In
practice PHA is used as a scenario-decomposition heuristic:
\begin{itemize}[leftmargin=1.6em,itemsep=2pt]
  \item It often finds good solutions for stochastic mixed-integer and
        nonconvex problems, especially when the scenario subproblems are
        easy and the coupling is weak \citep{watson2011}.
  \item It can stall at a non-consensus fixed point if $r$ is poorly chosen or
        the scenario subproblems are solved too loosely.
  \item Different initial policies and different penalty schedules can lead to
        different local solutions; use multi-start and compare the aggregated
        objective.
  \item Always validate the final consensus policy on the full SAA objective
        and on an independent test archive; do not trust the PHA objective
        alone.
\end{itemize}

\paragraph{Penalty selection.}
The penalty $r$ controls the trade-off between scenario freedom and consensus.
A small $r$ lets scenarios drift, so the consensus residual decays slowly; a
large $r$ forces consensus but makes the subproblems stiff and can slow the
inner solver.  Practical rules:
\begin{enumerate}[leftmargin=1.6em,itemsep=2pt]
  \item start with $r$ of the order of the curvature of $J$ with respect to
        $\qq$ (estimate it by a finite-difference Hessian-vector product on a
        few scenarios);
  \item increase $r$ by a factor of $2$--$10$ if the consensus residual has
        not decreased by a fixed factor over several iterations;
  \item decrease $r$ if the multipliers grow without bound or the subproblem
        solver fails;
  \item use scenario-dependent penalties $r_m$ when the scenario costs have
        very different scales, but keep the aggregation formula consistent.
\end{enumerate}

\paragraph{Inner-solver accuracy.}
Early PHA iterations do not need an accurate scenario solution.  Use a loose
inner tolerance (e.g.\ $10^{-2}$ relative) for the first few iterations and
tighten it as the consensus residual decreases.  Warm-start each subproblem
from the previous scenario policy.  This ``inexact PHA'' is much cheaper than
solving every subproblem to high accuracy.

\paragraph{Parallelism.}
The $M$ scenario subproblems in \eqref{eq:pha-sub} are independent.  This is
PHA's main advantage: on a cluster, the wall-clock cost per iteration is the
cost of one (or a few) scenario solves, not $M$ sequential solves.  The
aggregation and multiplier update are negligible.  If $M$ is larger than the
number of workers, use scenario bundling: group scenarios into bundles, solve
one subproblem per bundle with the bundle average objective, and treat the
bundle as one scenario in PHA.  This reduces communication and can improve
stability.

\subsection{PHA with chance constraints}
\label{sec:pha-chance}

The chance constraint couples scenarios, so it does not fit the separable form
\eqref{eq:consensus-program} directly.  There are four practical ways to
combine PHA with the reliability requirement.  They differ in how much
conservatism they introduce and in whether the resulting subproblems remain
scenario-separable.

\subsubsection{CVaR surrogate with an auxiliary variable}

Replace the chance constraint by the sufficient CVaR condition
\begin{equation}
\label{eq:pha-cvar}
\CVaR_{1-\alpha}(G(\qq,\omega))
=\min_{\eta\in\R}\left\{
\eta+\frac{1}{\alpha}\E\bigl[(G(\qq,\omega)-\eta)_+\bigr]
\right\}\le0.
\end{equation}
For the SAA, introduce the common variable $\eta$ and write
\begin{equation}
\label{eq:pha-cvar-saa}
\min_{\qq,\eta}\ \sum_{m=1}^{M}\pi_m J_m(\qq)
\quad\text{s.t.}\quad
\eta+\frac{1}{\alpha}\sum_{m=1}^{M}\pi_m
\bigl(G_m(\qq)-\eta\bigr)_+\le0.
\end{equation}
The positive-part term is scenario-separable for fixed $\eta$, but the
constraint is not.  Dualize the CVaR constraint with a multiplier
$\lambda\ge0$:
\begin{equation}
\label{eq:pha-cvar-lag}
L(\qq,\eta,\lambda)
=\sum_{m=1}^{M}\pi_m\left[
J_m(\qq)+\frac{\lambda}{\alpha}
\bigl(G_m(\qq)-\eta\bigr)_+
\right]
+\lambda\eta.
\end{equation}
For fixed $\lambda$, the problem is separable across scenarios given the common
$\eta$.  Treat $\eta$ as an additional consensus variable in PHA: introduce
copies $\eta^m$, add the consensus constraint $\eta^m=\eta$, and augment
\eqref{eq:pha-cvar-lag} with the corresponding multiplier and proximal term.
The scenario subproblem becomes
\begin{equation}
\label{eq:pha-cvar-sub}
\min_{\qq^m,\eta^m}\ 
\pi_m J_m(\qq^m)
+\frac{\lambda\pi_m}{\alpha}\bigl(G_m(\qq^m)-\eta^m\bigr)_+
+\text{(consensus terms for $\qq^m$ and $\eta^m$)}.
\end{equation}
After the PHA sweep, update $\lambda$ by projected subgradient ascent on the
CVaR constraint:
\begin{equation}
\label{eq:pha-lambda}
\lambda^{k+1}
=\max\left(0,\
\lambda^k+s_k\left[
\eta^{k+1}+\frac{1}{\alpha}\sum_{m=1}^{M}\pi_m
\bigl(G_m(\qq^{k+1,m})-\eta^{k+1}\bigr)_+
\right]\right),
\end{equation}
with $s_k>0$ a decreasing step.  This is a PHA--Lagrangian hybrid: PHA handles
the nonanticipativity of $\qq$ and $\eta$, and the outer multiplier handles the
CVaR inequality.  The CVaR condition is conservative relative to the chance
constraint, so a feasible CVaR policy is reliable; the price is conservatism.

\subsubsection{Lagrangian relaxation of the empirical probability}

A more direct approach keeps the empirical probability.  Write the chance
constraint as
\begin{equation}
\label{eq:pha-prob}
(1-\alpha)-\frac{1}{M}\sum_{m=1}^{M}\ind\{G_m(\qq)\le0\}\le0.
\end{equation}
Dualize it with $\lambda\ge0$:
\begin{equation}
\label{eq:pha-prob-lag}
L(\qq,\lambda)
=\E[J(\qq,\omega)]
+\lambda\left[(1-\alpha)-\Prob(G(\qq,\omega)\le0)\right]
=\E\bigl[J(\qq,\omega)-\lambda\ind\{G(\qq,\omega)\le0\}\bigr]
+\lambda(1-\alpha).
\end{equation}
For fixed $\lambda$, the scenario objective is
\begin{equation}
\label{eq:pha-prob-scenario}
f_m^\lambda(\qq^m)=J_m(\qq^m)-\lambda\ind\{G_m(\qq^m)\le0\}.
\end{equation}
This is scenario-separable, so PHA applies directly with $f_m=f_m^\lambda$.
The outer multiplier is updated by
\begin{equation}
\label{eq:pha-prob-lambda}
\lambda^{k+1}
=\max\left(0,\
\lambda^k+s_k\left[(1-\alpha)
-\frac{1}{M}\sum_{m=1}^{M}\ind\{G_m(\qq^{k+1,m})\le0\}\right]\right).
\end{equation}
If the empirical compliance probability is below $1-\alpha$, the bracket is
positive, $\lambda$ increases, and the reward $-\lambda$ for compliant
scenarios becomes stronger.  This is a clean decomposition, but it has two
caveats:
\begin{itemize}[leftmargin=1.6em,itemsep=2pt]
  \item the indicator is nonsmooth, so the scenario subproblem should be
        solved by a derivative-free method or by a smooth sigmoid
        approximation;
  \item for nonconvex problems there may be a positive duality gap, so a
        zero of the Lagrangian dual does not guarantee primal feasibility.
        Validate the final policy out of sample and increase $\lambda$ or add
        a penalty if the reliability target is missed.
\end{itemize}

\subsubsection{Sigmoid smoothing and continuation}

Replace the indicator by the smooth approximation
\begin{equation}
\label{eq:pha-sigmoid}
\ind\{G_m(\qq)\le0\}
\approx
\sigma\!\left(-\frac{G_m(\qq)}{\varepsilon_s}\right),
\end{equation}
and use the same Lagrangian decomposition with
\begin{equation}
\label{eq:pha-sigmoid-scenario}
f_m^{\lambda,\varepsilon_s}(\qq^m)
=J_m(\qq^m)
-\lambda\,\sigma\!\left(-\frac{G_m(\qq^m)}{\varepsilon_s}\right).
\end{equation}
The scenario subproblem is now differentiable, so an adjoint gradient or a
surrogate gradient can be used.  Use a continuation schedule
$\varepsilon_s=10^{-1},10^{-2},10^{-3}$ and, for each $\varepsilon_s$, run
PHA and the outer $\lambda$ update to convergence (or a fixed budget).  Warm-
start the next stage from the previous consensus policy.  The sigmoid is
smooth but not exact: at the final stage, check the empirical probability with
the exact indicator on the test archive.

\subsubsection{Mixed-integer reliability formulation}

If the chance constraint must be represented exactly, introduce
$z_m\in\{0,1\}$ with $z_m=1$ meaning scenario $m$ is compliant:
\begin{equation}
\label{eq:pha-minlp}
\sum_{m=1}^{M}z_m\ge(1-\alpha)M,
\qquad
G_m(\qq)\le M_G(1-z_m),\quad m=1,\dots,M,
\end{equation}
where $M_G$ is a valid upper bound on $G$.  The $z_m$ are scenario-local; the
reliability budget is the only coupling.  Dualize the budget with
$\lambda\ge0$:
\begin{equation}
\label{eq:pha-minlp-lag}
L(\qq,\{z_m\},\lambda)
=\sum_{m=1}^{M}\pi_m\left[
J_m(\qq)+\lambda\bigl((1-\alpha)-z_m\bigr)
\right],
\end{equation}
subject to $G_m(\qq)\le M_G(1-z_m)$ and $z_m\in\{0,1\}$.  For fixed $\lambda$,
each scenario subproblem chooses $z_m\in\{0,1\}$ and $\qq^m$:
\begin{itemize}[leftmargin=1.6em,itemsep=2pt]
  \item if $z_m=1$, the scenario must satisfy $G_m(\qq^m)\le0$ and pays
        $J_m(\qq^m)-\lambda$;
  \item if $z_m=0$, the compliance constraint is inactive and the scenario pays
        $J_m(\qq^m)$.
\end{itemize}
PHA handles the consensus on $\qq$; the outer $\lambda$ update is
\begin{equation}
\label{eq:pha-minlp-lambda}
\lambda^{k+1}
=\max\left(0,\
\lambda^k+s_k\left[(1-\alpha)
-\frac{1}{M}\sum_{m=1}^{M}z_m^{k+1}\right]\right).
\end{equation}
This is exact in the sense that, for a given $\lambda$, the scenario problems
are the Lagrangian subproblems of the MINLP; but the dual may have a gap, and
the subproblems contain binaries.  It is practical only for small $M$ and when
the scenario subproblems are cheap or surrogate-based.  It is useful as a
reference against which the CVaR and sigmoid approximations are compared.

\subsection{PHA with PDE subproblems}

Each PHA scenario subproblem is a single-scenario PDE-constrained optimization
with a proximal term.  There are three ways to solve it:
\begin{enumerate}[leftmargin=1.6em,itemsep=2pt]
  \item \textbf{Adjoint + L-BFGS-B.}  The proximal term is cheap; the
        gradient of $f_m$ comes from the discrete adjoint.  This is the
        highest-fidelity option and the natural choice when the forward solver
        is differentiable.
  \item \textbf{Surrogate subproblem.}  Replace the PDE solve by the FNO or a
        POD reduced model, solve the subproblem with autodiff, and validate the
        resulting consensus policy with the high-fidelity model.  This is the
        practical choice when $M$ is large and each forward solve is
        expensive.  The surrogate must be trained on the rate range visited by
        PHA, and the final policy must be re-evaluated with the FV model.
  \item \textbf{Derivative-free subproblem.}  If the chance constraint is kept
        in exact indicator form and $N_q$ is small, use pattern search or
        CMA-ES for the scenario subproblem.  This is robust but expensive.
\end{enumerate}

A useful large-scale pattern is \emph{surrogate-PHA}: train the FNO on a
scenario/rate design, run PHA entirely on the surrogate to obtain a candidate
consensus policy and a set of active scenarios, then evaluate that policy and
the active scenarios with the high-fidelity model.  If the high-fidelity
reliability or cost differs from the surrogate prediction, add the active
scenarios to the training set, retrain, and repeat.  This combines PHA's
scenario parallelism with the surrogate's speed.

\subsection{Strengths and weaknesses of PHA}

\paragraph{Strengths.}
\begin{itemize}[leftmargin=1.6em,itemsep=2pt]
  \item Scenario decomposition: the $M$ subproblems are independent and
        parallelizable.
  \item Works with nonsmooth scenario objectives and constraints; the
        proximal term regularizes the subproblems.
  \item Naturally handles a large number of scenarios and scenario bundles.
  \item Provides scenario policies $\qq^{k,m}$ and multipliers $\ww^{k,m}$
        that diagnose which scenarios disagree with the consensus.
  \item With CVaR or a Lagrangian chance constraint, it extends to reliability
        requirements without giving up scenario decomposition.
\end{itemize}

\paragraph{Weaknesses.}
\begin{itemize}[leftmargin=1.6em,itemsep=2pt]
  \item No global-convergence guarantee for the nonconvex PDE problem; it is a
        heuristic in that regime.
  \item The penalty $r$ and the outer multiplier step $s_k$ need tuning.
  \item Each iteration requires $M$ scenario subproblem solves; if the
        subproblems are solved accurately and $M$ is large, the total cost can
        exceed a well-tuned deterministic-equivalent NLP.
  \item The chance-constraint coupling requires a dual/penalty layer, and the
        dual may have a gap.
  \item The consensus policy is an average of scenario policies; it may not be
        optimal for any single scenario, and its out-of-sample performance must
        be checked.
\end{itemize}

\paragraph{When to use PHA.}
Use PHA when (i) $M$ is large or expected to grow, (ii) the scenario
subproblems are parallelizable and cheap relative to communication, (iii) the
objective and state equations are scenario-separable, and (iv) a convex or
weakly coupled structure makes the augmented Lagrangian well behaved.  For a
small $M$ and a smooth problem with a reliable adjoint, a direct SAA NLP is
usually simpler and faster.  PHA is not a substitute for a good scenario
generator, a validated forward model, or an out-of-sample certificate.

% ======================================================================
\section{Other decomposition methods}
\label{sec:decomp}

\subsection{L-shaped method and Benders decomposition}

For a two-stage stochastic linear program with recourse, the L-shaped method
\citep{vanslyke1969} is the canonical decomposition.  The master problem
chooses the here-and-now decision and a lower bound on the recourse cost; the
subproblems, one per scenario, generate optimality cuts.  Benders decomposition
\citep{benders1962} is the same idea in a more general setting.  The extension
to the remediation problem is not direct because the recourse (the PDE state)
is nonlinear and the chance constraint introduces binaries or nonsmooth
indicators.  Useful variants are:
\begin{itemize}[leftmargin=1.6em,itemsep=2pt]
  \item \textbf{Generalized Benders} for smooth nonlinear recourse, with cuts
        from the Lagrangian dual of the scenario subproblem;
  \item \textbf{Integer L-shaped} or branch-and-cut for the MINLP chance
        constraint, with the reliability budget as the complicating constraint;
  \item \textbf{Linearized recourse} around a reference policy, used as a
        lower-bound generator inside a trust-region or branch-and-bound
        scheme.
\end{itemize}
These methods give bounds and finite convergence in the convex case, which is
valuable for verification.  They are less attractive as a first implementation
because the PDE subproblems are nonlinear and the chance constraint is not
naturally linear.

\subsection{Dual decomposition and Lagrangian relaxation}

Dual decomposition relaxes the coupling constraints with multipliers and solves
independent subproblems.  It is the same principle as the Lagrangian layer in
Section~\ref{sec:pha-chance}.  It is most useful when the coupling is a small
number of linear or convex constraints (capacity, budget, reliability) and the
scenario subproblems are easy.  The disadvantages are the usual ones:
duality gap for nonconvex problems, slow subgradient convergence, and the need
for a primal recovery step.  For the remediation problem, dual decomposition is
a building block inside PHA or a stochastic-gradient method, not a stand-alone
solution.

\subsection{Scenario reduction and bundling}

For very large $M$, reduce the scenario archive before optimizing:
\begin{itemize}[leftmargin=1.6em,itemsep=2pt]
  \item fast forward selection or k-means in the space of KL coefficients;
  \item Wasserstein-based scenario reduction;
  \item bundling for PHA (each bundle is one subproblem).
\end{itemize}
Scenario reduction is safe for the expected objective when the reduced
distribution is close in an appropriate Wasserstein metric.  It is more
dangerous for a chance constraint: the tail scenarios that determine
reliability may be exactly the ones that reduction discards.  Keep the tail,
or use importance sampling, and always validate on an independent
unreduced archive.

% ======================================================================
\section{Stochastic gradient and stochastic approximation}
\label{sec:sg}

\subsection{Expected-value objective}

If the chance constraint is replaced by a smooth penalty or a CVaR term, the
problem has the form
\begin{equation}
\label{eq:sg-problem}
\min_{\qq\in\mathcal{Q}_{\mathrm{ad}}}\
F(\qq)=\E\bigl[\phi(\qq,\omega)\bigr],
\end{equation}
where $\phi$ is a smooth (or smoothed) scenario cost.  The gradient is
\begin{equation}
\label{eq:sg-grad}
\nabla F(\qq)=\E\bigl[\nabla_{\qq}\phi(\qq,\omega)\bigr].
\end{equation}
A mini-batch of $B\ll M$ scenarios gives an unbiased estimate
\begin{equation}
\label{eq:sg-est}
g^k=\frac{1}{B}\sum_{b=1}^{B}\nabla_{\qq}\phi(\qq^k,\omega^{k,b}),
\end{equation}
and the projected stochastic gradient iteration is
\begin{equation}
\label{eq:sg-update}
\qq^{k+1}=\Proj_{\mathcal{Q}_{\mathrm{ad}}}
\bigl(\qq^k-\gamma_k g^k\bigr).
\end{equation}
Standard choices are $\gamma_k=\gamma_0/\sqrt{k}$ (robust, $O(1/\sqrt{K})$
for convex problems) or $\gamma_k=\gamma_0/k$ (slower tail, better final
accuracy), with Polyak--Ruppert averaging
$\bar\qq^K=\frac{1}{K}\sum_{k=1}^{K}\qq^k$.  Adaptive methods (AdaGrad, RMSProp,
Adam) are convenient when the components of $\qq$ have very different scales
or the gradient noise is heteroscedastic.  They do not remove the need to tune
the step size or to validate the final policy.

For nonconvex problems, the guarantee is stationarity of the averaged iterates
under decreasing steps and bounded variance, not global optimality.  The
practical advantage is that each iteration touches only $B$ scenarios, so the
method scales to very large $M$ and can be run online as new scenarios arrive.

\subsection{Constrained variants}

The chance constraint can be handled by:
\begin{itemize}[leftmargin=1.6em,itemsep=2pt]
  \item \textbf{sigmoid smoothing}: replace the indicator by
        $\sigma(-G/\varepsilon_s)$ and optimize
        $\E[J]+\rho[1-\alpha-\E\sigma(-G/\varepsilon_s)]_+^2$;
  \item \textbf{CVaR}: add $\lambda[\eta+\frac{1}{\alpha}\E(G-\eta)_+]$ and
        update $\eta$ and $\lambda$ in an inner/outer loop;
  \item \textbf{two-time-scale stochastic approximation}: update $\qq$ on a
        fast timescale and the multipliers $\lambda$ (and $\eta$) on a slow
        timescale;
  \item \textbf{penalty}: minimize $\E[J]+\rho[(1-\alpha)-\ind\{G\le0\}]_+^2$
        with a smoothed indicator and increase $\rho$ until feasibility.
\end{itemize}
The constrained stochastic approximation literature (e.g.\
\citealp{nemirovski2009}) gives convergence rates for convex and smooth
problems; for the nonconvex PDE problem, treat the rates as guidance rather
than a guarantee.

\subsection{Variance reduction}

Stochastic gradients are noisy; the following are essential in practice:
\begin{enumerate}[leftmargin=1.6em,itemsep=2pt]
  \item \textbf{Common random numbers (CRN).}  Use the same scenario batch for
        all candidates being compared; this reduces the variance of
        differences and of finite-difference gradients.
  \item \textbf{Antithetic variates.}  Pair $\xi$ with $-\xi$ in the KL
        coefficients when the functional is monotone.
  \item \textbf{Low-discrepancy sampling.}  Use Sobol or Latin hypercube points
        for the KL coefficients, especially when $N_{KL}$ is small.
  \item \textbf{Control variates.}  Use the FNO or a coarse-grid FV solution as
        a control variate for the fine-grid objective; the difference has much
        lower variance than either term alone.
  \item \textbf{Importance sampling.}  For rare compliance failures, tilt the
        scenario distribution toward high-concentration events and reweight.
        Check the estimator against plain Monte Carlo on a small subset.
\end{enumerate}

\subsection{Strengths and weaknesses}

\paragraph{Strengths.}  Scales to very large $M$; simple to implement; naturally
online; handles smooth penalties and CVaR; parallelizes over the batch.

\paragraph{Weaknesses.}  Noisy; step-size and penalty tuning; slow tail
convergence; constraint feasibility is not guaranteed without a careful
dual/penalty layer; a single run does not provide a reliability certificate.

% ======================================================================
\section{Derivative-free and global optimization}
\label{sec:df}

\subsection{Pattern search and MADS}

Pattern search (Hooke--Jeeves, generalized pattern search) evaluates the
objective on a set of poll directions around the current point and moves to a
better point; the step size shrinks when no improvement is found.  Mesh-
adaptive direct search (MADS) extends this to constrained problems with a
poll step and a search step, and has convergence theory for nonsmooth
functions under mild assumptions.  For the remediation problem:
\begin{itemize}[leftmargin=1.6em,itemsep=2pt]
  \item use the exact empirical chance constraint (no smoothing), with a
        filter or an augmented-Lagrangian constraint handling;
  \item use common random numbers so that differences between candidates are
        meaningful;
  \item start from a space-filling design and run several restarts;
  \item restrict to $N_q\lesssim5$--$10$ and a coarse SAA ($M\approx50$--$200$).
\end{itemize}
Pattern search is easy to implement, robust to noise, and a good way to find a
basin before switching to a gradient-based method.

\subsection{CMA-ES and differential evolution}

CMA-ES adapts a multivariate Gaussian search distribution to the objective
landscape; differential evolution maintains a population and uses mutation and
crossover.  Both are global-ish population methods that handle nonsmooth and
rugged objectives.  They are attractive for small $N_q$ and a coarse SAA, but
each candidate requires $M$ PDE solves, so the cost grows quickly.  Use:
\begin{itemize}[leftmargin=1.6em,itemsep=2pt]
  \item a small population and a coarse SAA for exploration;
  \item CRN across the population;
  \item a penalty or a constrained ranking for the chance constraint;
  \item a final local refinement with a gradient method or a trust-region
        surrogate.
\end{itemize}
Neither method provides a global optimality certificate for this problem.

\subsection{Bayesian optimization}

Bayesian optimization builds a Gaussian-process surrogate of the objective and
constraint functions and chooses the next evaluation by maximizing an
acquisition function (expected improvement, upper confidence bound, or a
constrained variant).  It is sample-efficient and well suited to very
expensive forward solves with $N_q\lesssim10$.  The difficulties are:
\begin{itemize}[leftmargin=1.6em,itemsep=2pt]
  \item the chance constraint is a probability, not a single smooth function;
        use a surrogate for $G_m$ or for the empirical probability, and a
        constrained acquisition function;
  \item the objective is noisy because of the Monte Carlo scenarios; use a
        noise-aware GP (nugget) and CRN;
  \item scaling to $N_q>20$ is poor;
  \item the GP is only a local model; validate the recommended policy with the
        high-fidelity model.
\end{itemize}

\subsection{Strengths and weaknesses}

\paragraph{Strengths.}  No gradient required; handles nonsmooth chance
constraints; robust to simulation noise; easy to parallelize the objective
evaluations; good for small $N_q$ and expensive FV.

\paragraph{Weaknesses.}  Cost grows with $N_q$ and with the number of
constraints; no global guarantee (except specialized global-search variants);
slow tail convergence; not competitive for $N_q$ in the hundreds.

% ======================================================================
\section{Surrogate and reduced-order optimization}
\label{sec:surrogate}

\subsection{Surrogate models}

A surrogate replaces the high-fidelity map
$(\qq,\omega)\mapsto(J,G,H)$ or $(\qq,\omega)\mapsto C(\cdot)$ by a cheap
approximation:
\begin{itemize}[leftmargin=1.6em,itemsep=2pt]
  \item \textbf{Fourier neural operator (FNO)}: learns the map from the
        conductivity field and the rate vector to the concentration field;
        differentiable, fast, and accurate for smooth parametric PDEs.
  \item \textbf{POD--Galerkin or POD--DEIM}: projects the state onto a
        reduced basis built from high-fidelity snapshots; fast for a fixed
        parameter range, less flexible for strong heterogeneity.
  \item \textbf{Gaussian process / polynomial chaos}: good for low-dimensional
        parameterizations and small data; uncertainty quantification is built
        in.
\end{itemize}
The surrogate is trained on a design of experiments over scenarios and rates.
For the remediation problem, the design must cover the no-pumping case, the
capacity case, asymmetric pumping patterns, and the compliance locations.
The training loss should include the field error, the well-concentration
error, and the mass-balance residual.

\subsection{Trust-region model management}

Optimize the surrogate inside a trust region around the current best policy:
\begin{equation}
\label{eq:tr}
\qq^{k+1}
=\argmin_{\qq\in\mathcal{Q}_{\mathrm{ad}}\cap\mathcal{B}(\qq^k,\Delta_k)}
\ \hat J_M(\qq)
\quad\text{s.t.}\quad
\hat c_M(\qq)\le0,
\end{equation}
where $\hat J_M$ and $\hat c_M$ are surrogate-based SAA functionals and
$\mathcal{B}(\qq^k,\Delta_k)$ is a trust region.  Evaluate the candidate with
the high-fidelity model on a validation set.  If the high-fidelity agreement
is good, accept the step and possibly expand $\Delta_k$; otherwise reject the
step, shrink $\Delta_k$, and add the candidate to the training set.  This is
the standard trust-region framework for simulation-based optimization and it
provides a practical safeguard against surrogate bias.

\subsection{Active learning and decision-focused validation}

A small field error does not imply a small decision error.  After solving the
surrogate problem, evaluate the recommended policy on an independent
high-fidelity test archive and compare:
\begin{itemize}[leftmargin=1.6em,itemsep=2pt]
  \item the SAA objective and the out-of-sample objective;
  \item the empirical compliance probability and the CVaR;
  \item the cleanup time and the residual mass;
  \item the active scenarios and the constraint margin.
\end{itemize}
If the discrepancy is too large, add the worst scenarios and the neighborhood
of the candidate policy to the training set, retrain, and repeat.  This
active-learning loop is the practical answer to surrogate bias: the surrogate
is not trusted until it has been tested on the decisions it recommends.

\subsection{Surrogate inside PHA and SAA}

The surrogate can be used at three levels:
\begin{enumerate}[leftmargin=1.6em,itemsep=2pt]
  \item \textbf{Inner level}: replace the PDE solve in each PHA scenario
        subproblem or each SAA objective evaluation;
  \item \textbf{Outer level}: solve the entire SAA/PHA problem on the
        surrogate, then validate the consensus policy with the FV model;
  \item \textbf{Control-variate level}: use the surrogate as a control variate
        to reduce the variance of the high-fidelity SAA estimator.
\end{enumerate}
The inner level is the most aggressive and the most dependent on surrogate
accuracy; the control-variate level is the safest and often gives a large
variance reduction with a small bias correction.

\subsection{Strengths and weaknesses}

\paragraph{Strengths.}  Many scenario evaluations per second; differentiable
with respect to $\qq$; enables PHA and SAA at scales that are impossible with
the FV model; can be used for real-time decision support.

\paragraph{Weaknesses.}  Surrogate bias; training cost; extrapolation outside
the training range; the decision error can be much larger than the field error.
Always pair the surrogate with a high-fidelity validation and an active-
learning loop.

% ======================================================================
\section{Chance-constraint-specific algorithms}
\label{sec:chance}

This section collects the reformulations that can be plugged into any of the
outer algorithms above.

\subsection{Exact mixed-integer reformulation}

With binaries $z_m\in\{0,1\}$,
\begin{equation}
\label{eq:minlp}
\sum_{m=1}^{M}z_m\ge(1-\alpha)M,
\qquad
G_m(\qq)\le M_G(1-z_m),\quad m=1,\dots,M.
\end{equation}
This is exact but is a mixed-integer nonlinear program with $M$ binaries and
$M$ PDE state constraints.  It is useful only for tiny $N_q$ and moderate $M$,
and even then the PDE evaluations dominate.  It is the reference against which
the CVaR and sigmoid approximations are compared.  The Lagrangian version in
Section~\ref{sec:pha-chance} is the scenario-decomposed way to use it.

\subsection{CVaR reformulation}

The CVaR constraint
\begin{equation}
\label{eq:cvar}
\CVaR_{1-\alpha}(G)
=\min_{\eta\in\R}\left\{
\eta+\frac{1}{\alpha}\E\bigl[(G-\eta)_+\bigr]
\right\}\le0
\end{equation}
is a sufficient condition for the chance constraint and is convex in the
probability space.  For the SAA, introduce $\eta$ and solve
\begin{equation}
\label{eq:cvar-saa}
\min_{\qq,\eta}\ J_M(\qq)
\quad\text{s.t.}\quad
\eta+\frac{1}{\alpha M}\sum_{m=1}^{M}
\bigl(G_m(\qq)-\eta\bigr)_+\le0.
\end{equation}
The positive part is nonsmooth at $G_m=\eta$; smooth it by
$(\cdot)_+\approx\frac{1}{2}(\cdot+\sqrt{(\cdot)^2+\epsilon^2})$ or use a
nonsmooth solver.  CVaR is more conservative than the chance constraint: a
CVaR-feasible policy is reliable, but a chance-feasible policy may be rejected.
Use CVaR when the tail of the violation distribution matters and when a
convex-in-probability surrogate is desirable.

\subsection{Sigmoid smoothing}

The empirical probability is approximated by
\begin{equation}
\label{eq:sigmoid}
\Prob_M(G\le0)
\approx\frac{1}{M}\sum_{m=1}^{M}
\sigma\!\left(-\frac{G_m(\qq)}{\varepsilon_s}\right).
\end{equation}
As $\varepsilon_s\downarrow0$, the approximation converges to the empirical
probability for almost every $\qq$.  In practice, use a continuation schedule
$\varepsilon_s=10^{-1},10^{-2},10^{-3}$ and re-optimize at each stage.  The
sigmoid is easy to differentiate and works well with adjoint or surrogate
gradients.  The drawback is that a small $\varepsilon_s$ produces very small
gradients when all scenarios are far from the boundary; a moderate
$\varepsilon_s$ is used during early iterations.  The final policy must be
checked with the exact indicator on an independent test archive.

\subsection{Difference-of-convex and successive convex approximation}

The empirical probability constraint can be written as a difference of convex
functions in some formulations, or approximated by a sequence of convex
surrogates (successive convex approximation, SCA).  The idea is to linearize
or convexify the nonconvex part around the current policy, solve a convex
subproblem, and use a trust region or line search.  This is attractive when
the dynamics can be linearized and the chance constraint is the only source of
nonconvexity.  For the full PDE problem, the convex subproblem is still a
PDE-constrained problem, so the practical benefit is mainly in the constraint
handling, not in the PDE cost.

\subsection{Penalty and augmented-Lagrangian methods}

A generic penalty formulation is
\begin{equation}
\label{eq:penalty}
\min_{\qq\in\mathcal{Q}_{\mathrm{ad}}}\
J_M(\qq)
+\rho\left[\left(1-\alpha-\Prob_M(G\le0)\right)_+\right]^2
+\rho_h\left[\left(1-\alpha_h-\Prob_M(H\ge0)\right)_+\right]^2.
\end{equation}
Use a smoothed probability and increase $\rho,\rho_h$ until the constraints
are satisfied at the candidate solution.  This is simple and works with any
outer solver, but the penalty parameters can be large and the final
feasibility must be checked independently.  The augmented Lagrangian is
preferable when accurate constraint gradients are available, because it avoids
some of the ill-conditioning of a pure quadratic penalty.

\subsection{Sample size and validation}

The empirical probability $\hat p_M$ has standard error
$\sqrt{\hat p_M(1-\hat p_M)/M}$.  To resolve a violation probability
$\alpha$, one needs $M\gg1/\alpha$; for $\alpha=0.05$, $M=100$ is too small
and $M=1000$--$5000$ is typical.  The Dvoretzky--Kiefer--Wolfowitz inequality
gives the distribution-free bound
\begin{equation}
\label{eq:dkw}
\Prob\left(\sup_{\qq}\abs{\hat p_M(\qq)-p(\qq)}>\varepsilon\right)
\le2\exp(-2M\varepsilon^2),
\end{equation}
so $M\ge\frac{1}{2\varepsilon^2}\ln\frac{2}{\delta}$ for accuracy
$\varepsilon$ with confidence $1-\delta$.  For $\varepsilon=0.01$ and
$\delta=0.05$, this is $M\approx1.85\times10^4$; this is the sample size for a
uniform probability certificate, not the sample size needed to find a good
policy.  The practical recipe is: $M=100$--$500$ for exploration, increase
until the out-of-sample objective and reliability stabilize, and use
$M=2000$--$5000$ (or more) for the final certificate.

% ======================================================================
\section{Risk-averse and distributionally robust extensions}
\label{sec:risk}

If the expected cost is replaced by a mean--CVaR objective,
\begin{equation}
\label{eq:meancvar}
(1-\beta)\E[J(\qq,\omega)]
+\beta\,\CVaR_{1-\alpha}\bigl(J(\qq,\omega)\bigr),
\end{equation}
the SAA problem remains convex in the probability variable and can be solved
by the same SAA/NLP or PHA machinery with an auxiliary $\eta$.  The PHA
derivation is analogous to Section~\ref{sec:pha-chance}: treat $\eta$ as a
consensus variable and dualize or penalize the CVaR term.

Distributionally robust optimization (DRO) replaces the expectation by a
worst-case expectation over an ambiguity set $\mathcal{P}$:
\begin{equation}
\label{eq:dro}
\min_{\qq\in\mathcal{Q}_{\mathrm{ad}}}\
\sup_{\Prob'\in\mathcal{P}}\E_{\Prob'}[J(\qq,\omega)]
\quad\text{s.t.}\quad
\Prob'(G\le0)\ge1-\alpha\ \ \forall\Prob'\in\mathcal{P}.
\end{equation}
For a Wasserstein ball, the inner problem can sometimes be reformulated as a
finite-dimensional convex program or solved by cutting planes; for a
moment-ambiguity set, it becomes a semidefinite or second-order cone program.
DRO is more conservative than the SAA chance constraint and is best treated as
a sensitivity analysis after the baseline policy is certified.  It is not the
first algorithm to implement.

% ======================================================================
\section{Recommended hybrid roadmap}
\label{sec:hybrid}

The following staged approach matches the gate-driven philosophy of the
forward-model plans.  Each stage produces an artifact that can be checked
before the next stage is started.

\subsection{Stage 0: freeze the scenario and validation protocol}

\begin{enumerate}[leftmargin=1.6em,itemsep=2pt]
  \item Freeze the random-field model, $N_{KL}$, the well layout, the cost
        coefficients, and the chance-constraint parameters.
  \item Generate a training archive and a disjoint test archive with
        independent seeds; store the KL coefficients and all random variables.
  \item Precompute the unit-rate flow responses and the well operators.
  \item Check the forward model against the trusted 1D/2D/3D references; do not
        optimize until the forward gates are green.
\end{enumerate}

\subsection{Stage 1: coarse exploration}

Run a derivative-free method (pattern search or CMA-ES) on a coarse SAA
($M\approx50$--$200$) with the exact empirical chance constraint and common
random numbers.  Keep the best few candidates.  This stage is robust and gives
a good basin; it is not the final answer.

\subsection{Stage 2: SAA refinement with adjoint gradients}

For each candidate, solve the SAA problem with the adjoint gradient and a
smooth (CVaR or sigmoid) chance constraint on the full training archive.  Use
L-BFGS-B, SQP, or a trust-region method, with a continuation in the smoothing
parameter and a multi-start around the coarse candidates.  Check the adjoint
gradient against a finite difference on a smooth test problem before trusting
the refinement.

\subsection{Stage 3: large-scenario decomposition or surrogate}

If $M$ is large and the scenario subproblems are parallelizable, run PHA on
the full archive, with the chance constraint handled by the CVaR/Lagrangian
layer.  If the FV solve is the bottleneck, train or update the FNO, run
surrogate-PHA, and validate the consensus policy with the FV model.  Use
active learning to add the worst scenarios and the candidate neighborhood to
the training set.

\subsection{Stage 4: out-of-sample certificate}

Evaluate the final policy on the independent test archive and report:
\begin{align}
\hat J_{\mathrm{test}}
&=\frac{1}{M_{\mathrm{test}}}\sum_{m\in\mathcal{T}}J(\qq^\star,\omega^m),
\label{eq:Jtest}\\
\hat p_{\mathrm{test}}
&=\frac{1}{M_{\mathrm{test}}}\sum_{m\in\mathcal{T}}
\ind\{G(\qq^\star,\omega^m)\le0\},
\label{eq:ptest}\\
\widehat{\CVaR}_{\mathrm{test}}
&=\min_{\eta}\left\{\eta+\frac{1}{\alpha M_{\mathrm{test}}}
\sum_{m\in\mathcal{T}}\bigl(G(\qq^\star,\omega^m)-\eta\bigr)_+
\right\}.
\label{eq:cvartest}
\end{align}
Accept the policy only if $\hat p_{\mathrm{test}}\ge1-\alpha-\varepsilon_p$
and the test objective is within a stated tolerance of the training objective.
Repeat the SAA with $K$ independent seeds and report the best policy and the
spread; a single run is not a certificate.

\begin{algorithm}[htbp]
\caption{Hybrid algorithm with an optional PHA stage}
\label{alg:hybrid}
\begin{algorithmic}[1]
\Require random-field model, fixed wells, cost coefficients, $\alpha$,
         $M_{\mathrm{coarse}}$, $M_{\mathrm{fine}}$, tolerances
\Ensure policy $\qq^\star$, out-of-sample certificate, run manifest
\State Generate training and test archives; freeze seeds and $N_{KL}$
\State Precompute unit-rate flow responses; verify the forward model
\State \textbf{Stage 1:} pattern search / CMA-ES on a coarse SAA with the
       exact empirical chance constraint; keep the best candidates
\State \textbf{Stage 2:} adjoint-based SAA refinement with CVaR or sigmoid
       smoothing; multi-start from Stage 1
\If{$M$ is large and scenario subproblems are parallelizable}
  \State \textbf{Stage 3a:} run Rockafellar--Wets PHA with the
         CVaR/Lagrangian chance-constraint layer; aggregate the scenario
         policies
\EndIf
\If{a validated FNO/POD surrogate is available}
  \State \textbf{Stage 3b:} run surrogate-PHA or surrogate-SAA; validate the
         candidate with the FV model; add active scenarios and retrain if the
         decision error exceeds tolerance
\EndIf
\State \textbf{Stage 4:} evaluate $\qq^\star$ on the test archive; compute
       $\hat J_{\mathrm{test}}$, $\hat p_{\mathrm{test}}$,
       $\widehat{\CVaR}_{\mathrm{test}}$, and confidence intervals
\State Repeat the SAA with $K$ independent seeds; report the best policy and
       the spread
\State \Return $\qq^\star$, metrics, figures, and \texttt{run\_config.json}
\end{algorithmic}
\end{algorithm}

% ======================================================================
\section{Verification and diagnostics}
\label{sec:verification}

The following checks are the minimum for a trustworthy stochastic optimization
run.  They are cheap compared with the optimization itself and catch most of
the mistakes that produce an apparently good but unreliable policy.

\begin{enumerate}[leftmargin=1.6em,itemsep=2pt]
  \item \textbf{Adjoint gradient check.}  On a smooth test problem (one well,
        homogeneous $K$, no chance constraint), compare the adjoint gradient
        with a central finite difference for several random directions and
        step sizes.  A relative error below $10^{-4}$ in double precision is
        the acceptance gate.  If it fails, stop and fix the adjoint, the
        boundary conditions, or the well regularization.
  \item \textbf{SAA replication.}  Solve $K$ independent SAA problems with the
        same $M$ but different seeds.  Report the best objective, the mean,
        and the spread.  A large spread means $M$ is too small or the problem
        is too nonconvex for a single run.
  \item \textbf{Out-of-sample reliability.}  Evaluate the policy on an
        independent test archive.  The empirical compliance probability must
        meet $1-\alpha$ within the Monte Carlo error; report a confidence
        interval, not only a point estimate.
  \item \textbf{PHA consensus diagnostics.}  Plot the consensus residual
        $\max_m\norm{\qq^{k,m}-\qq^k}$, the multiplier norm, and the
        aggregated objective versus iteration.  A stalled residual with a
        growing multiplier norm indicates a bad penalty or an infeasible
        chance-constraint layer.
  \item \textbf{PHA vs deterministic equivalent.}  On a small convex toy
        problem (or a linearized version of the remediation problem), verify
        that PHA converges to the same objective as the direct SAA NLP.  This
        validates the implementation before using PHA on the nonconvex
        problem.
  \item \textbf{Surrogate decision error.}  Compare the surrogate-recommended
        policy with its high-fidelity evaluation on the decision-relevant
        metrics (cost, reliability, cleanup time, extracted mass), not only on
        the field error.
  \item \textbf{Sensitivity.}  Vary $M$, $N_{KL}$, the penalty $r$, the
        smoothing parameter $\varepsilon_s$, and the chance-constraint level
        $\alpha$.  Report how the policy and the reliability change.  A policy
        that is only good at one setting is not robust.
\end{enumerate}

% ======================================================================
\section{Summary and decision guide}
\label{sec:summary}

\begin{enumerate}[leftmargin=1.6em,itemsep=3pt]
  \item \textbf{Start with SAA + adjoint + L-BFGS-B/SQP.}  It is the simplest
        method that uses the structure of the problem and gives exact local
        gradients.  Use CVaR or sigmoid smoothing for the chance constraint,
        with a continuation in the smoothing parameter.
  \item \textbf{Use Rockafellar--Wets PHA when scenario decomposition is the
        bottleneck.}  It is the natural algorithm for large $M$, parallel
        scenario subproblems, and a separable expected-cost objective.  For
        chance constraints, combine it with a CVaR/Lagrangian layer or the
        mixed-integer reliability formulation, and validate the consensus
        policy out of sample.
  \item \textbf{Use derivative-free methods for small $N_q$ and nonsmooth
        chance constraints.}  Pattern search, CMA-ES, and Bayesian
        optimization are robust exploration tools; they are not scalable to
        many wells or many periods.
  \item \textbf{Use a surrogate when the FV solve is the bottleneck.}  The FNO
        or a POD model makes PHA and SAA affordable at large $M$; pair it with
        a trust region, active learning, and a high-fidelity certificate.
  \item \textbf{Use stochastic gradient methods for very large or online
        scenario streams.}  They scale well but need variance reduction and a
        careful constraint layer.
  \item \textbf{Always finish with an independent out-of-sample certificate.}
        The SAA objective is not a reliability guarantee; the test archive,
        multiple SAA replications, and the diagnostics in
        Section~\ref{sec:verification} are what make the result credible.
\end{enumerate}

The recommended default for the first end-to-end run is therefore:
\[
\boxed{
\text{coarse pattern search}
\;\longrightarrow\;
\text{adjoint SAA with CVaR/sigmoid}
\;\longrightarrow\;
\text{PHA or surrogate-PHA if }M\text{ is large}
\;\longrightarrow\;
\text{out-of-sample certificate.}
}
\]
This sequence uses each algorithm where it is strongest and keeps a
high-fidelity check at the end.

% ======================================================================
\appendix

\section{PHA derivation details}
\label{app:pha}

The augmented Lagrangian \eqref{eq:pha-L} is the standard augmented Lagrangian
for the constraints $\qq^m-\qq=0$ with multipliers $\pi_m\ww^m$ and penalty
$\pi_m r$.  The aggregate step \eqref{eq:pha-q} follows from
\begin{equation}
\label{eq:pha-deriv}
\nabla_{\qq}L_r
=-\sum_{m=1}^{M}\pi_m\ww^m
-r\sum_{m=1}^{M}\pi_m(\qq^m-\qq)=0,
\end{equation}
so
\begin{equation}
\label{eq:pha-deriv2}
\qq=\sum_{m=1}^{M}\pi_m\qq^m
+\frac{1}{r}\sum_{m=1}^{M}\pi_m\ww^m,
\end{equation}
and the projection onto $\mathcal{Q}_{\mathrm{ad}}$ gives
\eqref{eq:pha-aggregate}.  The multiplier update \eqref{eq:pha-w} is the
standard dual ascent step for the same augmented Lagrangian.  If the
multipliers are initialized at zero, the weighted sum of the multipliers
remains zero:
\begin{equation}
\label{eq:pha-mult-sum}
\sum_{m=1}^{M}\pi_m\ww^{k+1,m}
=\sum_{m=1}^{M}\pi_m\ww^{k,m}
+r\sum_{m=1}^{M}\pi_m(\qq^{k+1,m}-\qq^{k+1})
=\sum_{m=1}^{M}\pi_m\ww^{k,m},
\end{equation}
because $\qq^{k+1}=\sum_m\pi_m\qq^{k+1,m}$ when the projection is inactive.
Hence the aggregate is the weighted average \eqref{eq:pha-average}.

The PHA iteration is a Gauss--Seidel sweep on the augmented Lagrangian: first
minimize over the consensus $\qq$, then over the scenario copies $\qq^m$, then
update the multipliers.  The order can be reversed; the convergence theory is
the same under the standard assumptions.

\section{CVaR and chance-constraint equivalence}
\label{app:cvar}

For a random variable $Z$ and a confidence level $1-\alpha$, the CVaR is
\begin{equation}
\label{eq:app-cvar}
\CVaR_{1-\alpha}(Z)
=\min_{\eta\in\R}\left\{
\eta+\frac{1}{\alpha}\E\bigl[(Z-\eta)_+\bigr]
\right\}.
\end{equation}
The minimizer $\eta^\star$ is a $(1-\alpha)$-quantile of $Z$.  The key
inequality is
\begin{equation}
\label{eq:app-cvar-ineq}
\Prob(Z\le0)\ge1-\alpha
\quad\Longleftarrow\quad
\CVaR_{1-\alpha}(Z)\le0.
\end{equation}
Thus a CVaR-feasible policy satisfies the chance constraint, but the converse
is false: CVaR controls the average of the worst $\alpha$ fraction, not only
the probability of violation.  This is why the CVaR reformulation is
conservative and why it is useful as a tractable surrogate when the chance
constraint is nonsmooth.

For the SAA, the empirical CVaR is
\begin{equation}
\label{eq:app-cvar-saa}
\CVaR_{1-\alpha}^M(Z)
=\min_{\eta}\left\{
\eta+\frac{1}{\alpha M}\sum_{m=1}^{M}(Z_m-\eta)_+
\right\},
\end{equation}
and the constraint $\CVaR_{1-\alpha}^M(G(\qq))\le0$ is equivalent to the
existence of $\eta$ satisfying the inequality in \eqref{eq:cvar-saa}.  This is
the form used in the PHA--CVaR hybrid of Section~\ref{sec:pha-chance}.

\section{PHA convergence notes}
\label{app:convergence}

The convergence result in Proposition~\ref{prop:pha-convex} is the
Rockafellar--Wets theorem for convex scenario problems.  The essential
ingredients are:
\begin{enumerate}[leftmargin=1.6em,itemsep=2pt]
  \item the scenario objectives are proper, closed, convex functions;
  \item the feasible set is convex and closed;
  \item a constraint qualification holds for the consensus constraints;
  \item the augmented Lagrangian subproblems are solved exactly (or with
        summable error).
\end{enumerate}
Under these conditions, the sequence of aggregates converges to an optimal
solution and the multipliers converge to optimal dual variables.  For
polyhedral scenario problems, finite convergence is possible but not
guaranteed.

For nonconvex problems, the standard result is local: if the iteration
converges, it converges to a stationary point of the augmented Lagrangian
under a constraint qualification and a second-order condition.  There is no
general global guarantee, and the algorithm can cycle or stall.  Practical
remedies are adaptive penalties, multiple starts, inexact inner solves with a
summable accuracy schedule, and a final high-fidelity validation.

% ======================================================================
\bibliographystyle{plainnat}
\begin{thebibliography}{99}

\bibitem[Benders(1962)]{benders1962}
J.~F. Benders.
Partitioning procedures for solving mixed-variables programming problems.
\textit{Numerische Mathematik}, 4:238--252, 1962.

\bibitem[Birge and Louveaux(2011)]{birge2011}
J.~R. Birge and F.~Louveaux.
\textit{Introduction to Stochastic Programming}.
Springer, New York, 2nd edition, 2011.

\bibitem[Conn et~al.(2009)]{conn2009}
A.~R. Conn, K.~Scheinberg, and L.~N. Vicente.
\textit{Introduction to Derivative-Free Optimization}.
SIAM, Philadelphia, 2009.

\bibitem[Hansen(2016)]{hansen2016}
N.~Hansen.
The CMA evolution strategy: A tutorial.
arXiv:1604.00772, 2016.

\bibitem[Kleywegt et~al.(2002)]{kleywegt2002}
A.~J. Kleywegt, A.~Shapiro, and T.~Homem-de-Mello.
The sample average approximation method for stochastic discrete optimization.
\textit{SIAM Journal on Optimization}, 12(2):479--502, 2002.

\bibitem[Nemirovski et~al.(2009)]{nemirovski2009}
A.~Nemirovski, A.~Juditsky, G.~Lan, and A.~Shapiro.
Robust stochastic approximation approach to stochastic programming.
\textit{SIAM Journal on Optimization}, 19(4):1574--1609, 2009.

\bibitem[Nocedal and Wright(2006)]{nocedal2006}
J.~Nocedal and S.~J. Wright.
\textit{Numerical Optimization}.
Springer, New York, 2nd edition, 2006.

\bibitem[Peherstorfer et~al.(2018)]{peherstorfer2018}
B.~Peherstorfer, K.~Willcox, and M.~Gunzburger.
Survey of multifidelity methods in uncertainty propagation, inference, and
optimization.
\textit{SIAM Review}, 60(3):550--591, 2018.

\bibitem[Rockafellar and Uryasev(2000)]{rockafellar2000}
R.~T. Rockafellar and S.~Uryasev.
Optimization of conditional value-at-risk.
\textit{Journal of Risk}, 2(3):21--41, 2000.

\bibitem[Rockafellar and Wets(1991)]{rockafellar1991}
R.~T. Rockafellar and R.~J.-B. Wets.
Scenarios and policy aggregation in optimization under uncertainty.
\textit{Mathematics of Operations Research}, 16(1):119--147, 1991.

\bibitem[Shapiro et~al.(2009)]{shapiro2009}
A.~Shapiro, D.~Dentcheva, and A.~Ruszczy\'nski.
\textit{Lectures on Stochastic Programming: Modeling and Theory}.
SIAM, Philadelphia, 2009.

\bibitem[Shahriari et~al.(2016)]{shahriari2016}
B.~Shahriari, K.~Swersky, Z.~Wang, R.~P. Adams, and N.~de~Freitas.
Taking the human out of the loop: A review of Bayesian optimization.
\textit{Proceedings of the IEEE}, 104(1):148--175, 2016.

\bibitem[Van Slyke and Wets(1969)]{vanslyke1969}
R.~M. Van Slyke and R.~J.-B. Wets.
L-shaped linear programs with applications to optimal control and stochastic
programming.
\textit{SIAM Journal on Applied Mathematics}, 17(4):663--681, 1969.

\bibitem[Watson and Woodruff(2011)]{watson2011}
J.-P. Watson and D.~L. Woodruff.
Progressive hedging innovations for a class of stochastic mixed-integer
resource allocation problems.
\textit{Computational Management Science}, 8(4):355--370, 2011.

\end{thebibliography}

\end{document}
